\begin{helper}
Let \(\Omega\) be a measurable space, let \(V\) be a finite coordinate type,
let \(\nu\) be a measure on \(\Omega\), let \(F\subseteq\Omega\), and let
\(p:\Omega\to\mathbb R^V\) be measurable.  Let \(M\) be an
extended nonnegative real number.  Assume that \(\nu(F)\) is finite, that
\(M=\nu(F)\), and that the lower-rectangle law holds: for every
\(t:V\to\mathbb R\) with \(0\le t_v\le1\) for all \(v\),
\[
\nu\{\omega:F(\omega)\text{ and }0\le p(\omega)_v\le t_v
\text{ for every }v\in V\}
=M\,\operatorname{ofReal}\!\left(\prod_{v\in V}t_v\right).
\]
If \(C\subseteq\mathbb R^V\) is Borel and contained in the unit cube
\([0,1]^V\), then
\[
\nu\bigl(\{\omega:F(\omega),\ p(\omega)\in C\}
\cap\{\omega:p(\omega)_v\in[0,1]\text{ for every }v\in V\}\bigr)
= M\,\operatorname{vol}(C).
\]
\end{helper}
\begin{proof}
Let \(Q=[0,1]^V\) and let
\(Q_\Omega=\{\omega:p(\omega)_v\in[0,1]\text{ for every }v\in V\}\).
Restrict \(\nu\) to \(F\cap Q_\Omega\) and push it forward by \(p\); call
the resulting measure \(\mu\) on \(\mathbb R^V\).  Let
\(\rho=M\,\operatorname{vol}\!\restriction Q\).  Since \(\Omega\) has the
all-subsets measurable structure, \(p\) is measurable.

The hypothesis with \(t_v=1\) for every \(v\) gives
\[
\mu(\mathbb R^V)=\nu(F\cap Q_\Omega)=M.
\]
Also \(\operatorname{vol}(Q)=1\), by the finite-product volume formula for
\(\prod_v[0,1]\), and therefore \(\rho(\mathbb R^V)=M\).  Because
\(\mu(\mathbb R^V)\le\nu(F)<\infty\), \(\mu\) is finite; \(\rho\) has the
same finite total mass.

We next compare \(\mu\) and \(\rho\) on every product lower orthant
\(\{q:q_v\le s_v\text{ for all }v\}\), with arbitrary \(s:V\to\mathbb R\).
If some \(s_v<0\), then the orthant has empty intersection with \(Q\), so
both measures assign it mass \(0\).  If all \(s_v\ge0\), put
\(t_v=\min(s_v,1)\).  Then \(t_v\in[0,1]\), and intersecting the orthant
with \(Q\) is exactly the box \(\prod_v[0,t_v]\).  The lower-rectangle
hypothesis gives
\[
\mu\{q:q_v\le s_v\text{ for all }v\}
=M\,\operatorname{ofReal}\!\left(\prod_v t_v\right).
\]
The finite-product Lebesgue formula gives the same value for
\(\rho\), because \(\rho\) is \(M\) times Lebesgue measure restricted to
\(Q\).  Hence \(\mu\) and \(\rho\) agree on all finite product lower
orthants.

The one-dimensional lower rays generate the Borel sigma-algebra on
\(\mathbb R\), and finite products of these rays form a pi-system generating
the product Borel sigma-algebra on \(\mathbb R^V\).  Since \(\mu\) and
\(\rho\) are finite measures, agree on this generating pi-system, and have
the same total mass, the finite-measure uniqueness theorem gives
\(\mu=\rho\) on all Borel subsets of \(\mathbb R^V\).

Apply this equality to the Borel set \(C\).  The event in the statement is
precisely \(p^{-1}(C)\cap F\cap Q_\Omega\), whose \(\nu\)-measure is
\(\mu(C)\).  Since \(C\subseteq Q\),
\[
\rho(C)=M\,\operatorname{vol}(Q\cap C)=M\,\operatorname{vol}(C).
\]
Therefore the displayed event has measure \(M\,\operatorname{vol}(C)\), as
claimed.
\end{proof}
